% VI26 detailed challenge solutions -- VI/24-depth standard.

% VI26-DETAILED-CHALLENGE-01
\subsection*{Challenge 1 solution}
\begin{solution}
Let
\[
\Spec A=D(f_1)\cup\cdots\cup D(f_r),
\]
and suppose every \(A_{f_i}\) is Noetherian.  We show every ideal \(I\subseteq A\) is finitely
generated.

For each \(i\), the localized ideal \(I_{f_i}\) is finitely generated in the Noetherian ring
\(A_{f_i}\).  Choose finitely many elements of \(I\) whose localizations generate \(I_{f_i}\);
clearing denominators if necessary lets us choose actual elements of \(I\).  Let \(J\subseteq I\)
be generated by the union of all these finite sets.  Then
\[
I_{f_i}=J_{f_i}\qquad (1\le i\le r).
\]
Set \(M=I/J\).  Thus \(M_{f_i}=0\) for every \(i\).

Take \(m\in M\).  From \(M_{f_i}=0\) there is \(N_i\ge0\) with \(f_i^{N_i}m=0\).
Because the \(D(f_i)\) cover \(\Spec A\),
\[
V(f_1,\ldots,f_r)=\varnothing,
\]
so \((f_1,\ldots,f_r)=A\).  Replacing \(f_i\) by positive powers does not change the radical of the
generated ideal, hence
\[
(f_1^{N_1},\ldots,f_r^{N_r})=A.
\]
Choose \(a_i\in A\) with
\[
1=\sum_i a_i f_i^{N_i}.
\]
Then
\[
m=\sum_i a_i f_i^{N_i}m=0.
\]
Thus \(M=0\), so \(I=J\) is finitely generated.  Since \(I\) was arbitrary, \(A\) is Noetherian.

The proof isolates the local-global mechanism: finite generation is obtained locally, and the finite
distinguished cover converts local annihilation of \(I/J\) into global annihilation.
\end{solution}

% VI26-DETAILED-CHALLENGE-02
\subsection*{Challenge 2 solution}
\begin{solution}
Take
\[
A=k[u_1,u_2,u_3,\ldots]
\]
and in \(A[x]\) let
\[
I=(u_1,u_2,u_3,\ldots)A[x].
\]
Set
\[
B=A[x]/I.
\]
The algebra \(B\) is finite type over \(A\): the single element \(\bar x\) generates it.

We first prove that \(I\) is not finitely generated.  Suppose \(I=(h_1,\ldots,h_r)\).  Each
polynomial \(h_j\) has only finitely many coefficients, and each coefficient is a finite
\(A\)-linear combination of the \(u_i\).  Hence there is \(N\) such that every coefficient of every
\(h_j\) belongs to \((u_1,\ldots,u_N)\).  Therefore
\[
(h_1,\ldots,h_r)\subseteq (u_1,\ldots,u_N)A[x],
\]
which cannot contain the element \(u_{N+1}\in I\).  Contradiction.

It remains to justify that this prevents finite presentation even if one allows another finite set
of algebra generators.  We use the standard finite-generator invariance lemma:

\emph{Lemma.} If \(B\) is a finitely presented \(A\)-algebra and
\[
A[x_1,\ldots,x_n]\twoheadrightarrow B
\]
is any surjection with finitely many variables, then its kernel is finitely generated.

To see this, start from a finite presentation
\[
B=A[y_1,\ldots,y_m]/(r_1,\ldots,r_s).
\]
Express the images of the \(x_i\) as polynomials \(F_i(y)\), and the images of the \(y_j\) as
polynomials \(G_j(x)\).  Then the kernel in \(A[x_1,\ldots,x_n]\) is generated by the finitely many
substituted relations
\[
r_\ell(G(x))
\]
together with
\[
x_i-F_i(G(x)).
\]
For a kernel element \(h(x)\), write
\[
h(F(y))=\sum_{\ell=1}^s a_\ell(y)r_\ell(y).
\]
After substituting \(y=G(x)\), the expression \(h(F(G(x)))\) belongs to the ideal generated by
the \(r_\ell(G(x))\).  On the other hand,
\[
h(x)-h(F(G(x)))
\]
belongs to the ideal generated by the differences
\[
x_i-F_i(G(x)),
\]
by the polynomial-difference identity applied one variable at a time.  Hence every kernel element
belongs to the displayed finite set of generators.

Applying the lemma to the canonical surjection \(A[x]\twoheadrightarrow B\) would force \(I\) to be
finitely generated, contradiction.  Hence \(B\) is finite type but not finitely presented over
the non-Noetherian base \(A\).
\end{solution}

% VI26-DETAILED-CHALLENGE-03
\subsection*{Challenge 3 solution}
\begin{solution}
Consider a Cartesian square
\[
\begin{tikzcd}
X' \arrow[r] \arrow[d,"f'"'] & X \arrow[d,"f"]\\
S' \arrow[r] & S
\end{tikzcd}
\]
with \(f\) finite type.  We prove directly that \(f'\) is finite type.

Choose an affine cover \(V=\Spec A\subseteq S\).  Because \(f\) is finite type,
\(f^{-1}(V)\) has a finite affine cover
\[
U_j=\Spec B_j,\qquad j=1,\ldots,r,
\]
where each
\[
B_j=A[b_{j1},\ldots,b_{jn_j}].
\]
Now cover the open inverse image of \(V\) in \(S'\) by affine opens
\(V'=\Spec A'\).  Over such a \(V'\), the pullbacks of the \(U_j\) form the finite affine cover
\[
U'_j=U_j\times_VV'
=
\Spec(B_j\otimes_AA')
\]
of \(f'^{-1}(V')\).  Moreover
\[
B_j\otimes_AA'
=
A'[b_{j1}\otimes1,\ldots,b_{jn_j}\otimes1],
\]
so every \(U'_j\) is finite type over \(V'\).  We have therefore produced exactly the finite affine
cover required for \(f'\) to be finite type.

The abstract proof from VI/24–VI/26 is shorter: finite type is stable under arbitrary base change.
The direct proof explains why: tensor product carries a finite list of algebra generators to the
same finite list after scalar extension, while the finite affine cover remains finite.
\end{solution}

% VI26-DETAILED-CHALLENGE-04
\subsection*{Challenge 4 solution}
\begin{solution}
Let
\[
B=\mathbb Z[x,y]/(y^2-x^3-x),
\qquad
X=\Spec B.
\]
The ring \(B\) is generated by the two classes of \(x,y\) as a \(\mathbb Z\)-algebra, so the
structure morphism
\[
X\longrightarrow\Spec\mathbb Z
\]
is finite type.

The ring \(\mathbb Z\) is Noetherian.  Hilbert's basis theorem gives
\(\mathbb Z[x,y]\) Noetherian, and quotients of Noetherian rings are Noetherian.  Therefore \(B\)
is Noetherian and \(X\) is a Noetherian affine scheme.

For a point \(s\in\Spec\mathbb Z\), the fiber is obtained by tensoring with \(\kappa(s)\):
\[
X_s
=
\Spec\!\left(
\kappa(s)[x,y]/(y^2-x^3-x)
\right).
\]
This is generated by \(x,y\) over the field \(\kappa(s)\), so it is finite type.  For the generic
point the field is \(\mathbb Q\); for a closed point \((p)\) it is \(\mathbb F_p\).  No
factorization or smoothness analysis of the polynomial is needed: finite type of every fiber follows
formally from base-change stability.
\end{solution}

% VI26-DETAILED-CHALLENGE-05
\subsection*{Challenge 5 solution}
\begin{solution}
The basic implications are:
\[
\text{finite type}
\Longleftrightarrow
\text{locally finite type}+\text{quasi-compact}.
\]
Thus finite type always implies both local finite type and quasi-compactness, and the conjunction of
those two conditions always implies finite type.

Neither condition alone implies the other.  The morphism
\[
\coprod_{n\ge1}\Spec k\to\Spec k
\]
is locally finite type but not quasi-compact, hence not finite type.  On the other hand,
\[
\Spec k[x_1,x_2,\ldots]\to\Spec k
\]
is affine and therefore quasi-compact, but it is not locally finite type.

Now bring in Noetherianity.  If the base \(Y\) is locally Noetherian and \(X\to Y\) is locally
finite type, then \(X\) is locally Noetherian.  If \(Y\) is Noetherian and \(X\to Y\) is finite
type, then \(X\) is Noetherian.  These hypotheses on the base are essential.  For a non-Noetherian
ring \(A\), the identity
\[
\Spec A\to\Spec A
\]
is finite type, but its source need not be locally Noetherian.

Conversely, a Noetherian source does not force a morphism to be locally finite type, even over a
Noetherian field base.  Let
\[
K=k(t_1,t_2,t_3,\ldots)
\]
have infinite transcendence degree over \(k\).  Then \(\Spec K\) is a one-point Noetherian scheme,
but
\[
\Spec K\to\Spec k
\]
is not locally finite type because \(K\) is not a finitely generated \(k\)-algebra.

A Noetherian source does imply quasi-compactness of every morphism out of it: the inverse image of an
affine target open is an open subset of a Noetherian space, hence quasi-compact.  A merely locally
Noetherian source need not be quasi-compact; the infinite disjoint union of \(k\)-points is the
standard counterexample.

The useful classification is therefore:
\[
\begin{array}{c|c}
\text{Hypothesis} & \text{Valid conclusion}\\ \hline
\text{finite type} & \text{locally finite type and quasi-compact}\\
\text{lft + quasi-compact} & \text{finite type}\\
Y\text{ locally Noetherian + lft} & X\text{ locally Noetherian}\\
Y\text{ Noetherian + finite type} & X\text{ Noetherian}\\
X\text{ Noetherian} & \text{every }X\to Y\text{ is quasi-compact}
\end{array}
\]
All stronger converse implications listed in the challenge fail without additional hypotheses.
\end{solution}
